% Redrawn 2026-10-15; main labels 8 pt, notes 7.5 pt; actual fonts embedded in PDF.
% export-reference-width-mm: 112
% Original HER data-flow diagram. No relabeling edge creates real interactions.
\begin{tikzpicture}[x=\linewidth,y=1mm,
 font=\figurefont\fontsize{8}{10}\selectfont,
 box/.style={rounded corners=1.8pt,draw=FigureStroke,fill=FigureBlueFill,line width=.6pt,
   text width=23mm,minimum height=11mm,align=center,inner sep=1.5mm},
 flow/.style={-{Latex[length=2mm]},draw=FigureStroke,line width=.75pt}]
 \node[box] (task) at (.14,0) {任务生成\\目标$g$};
 \node[box] (env) at (.50,0) {真实交互\\$\beta_g,T$};
 \node[box] (trajectory) at (.86,0) {转移记录\\$(s_t,a_t,s_{t+1})$};
 \node[box,fill=FigureGreenFill] (relabel) at (.50,-26) {目标重标记\\$g'=(3,1)$};
 \node[box] (samples) at (.14,-26) {训练样本\\重算$r_t,\delta_t$};
 \draw[flow] (task) -- node[above] {采样} (env);
 \draw[flow] (env) -- node[above] {记录} (trajectory);
 \draw[flow] (trajectory.south) |- node[pos=.7,above] {已达位置} (relabel.east);
 \draw[flow] (relabel) -- (samples);
 \node[anchor=north,align=center,text=BookMuted,font=\FigureNoteFont\fontsize{7.5}{9.5}\selectfont] at (.5,-38)
   {上行：取得新交互记录\qquad 下行：重新解释已有经验};
\end{tikzpicture}
